% Statement and cited-input packet. Target proof is separately embargoed.
% Frozen II source SHA256: 74497bfe327fae25c9cdcb17334b3bf42a9793a9df7c0a4404d97568f2f760f0
% Target owner: cdp:st:D.1.13; tag: 01FY

\begin{proposition}[Refuted printed surjectivity and positivity, with a second example]\hypertarget{tag-01FY}{}\label{cdp:st:D.1.13}



For Wang's threefold $Y=X_2(\lambda)$ in \ref{cdp:st:D.1.6} one has
\[
K_Y\simeq f^*\OO_{\PP^1}(-1),\qquad
R^1f_*\OO_Y\simeq\OO\oplus\OO(-1),\qquad
R^2f_*\OO_Y\simeq\OO(-1).
\]
Consequently $h^q(Y,\OO_Y)=(1,1,0,0)$, and restriction to every
smooth two-torus fibre has rank one. This threefold satisfies
$H^1(Y;\Z)=H^2(Y;\Z)=0$, $a(Y)=1$, and is a second
counterexample to CDP20 Corollary 5.2 and Proposition 5.5.
Its integral homology torsion distinguishes it from the retained
integral homology sphere.
\end{proposition}

% Cited input cdp:st:D.1.12; statement or labelled formula only.
\begin{proposition}\hypertarget{tag-01FX}{}\label{cdp:st:D.1.12}


For a proper connected-fibre holomorphic map $f:Y\to C$ from a
compact complex threefold with $b_2(Y)=0$ and general fibre a
two-torus or a primary Kodaira surface, the cup-product homomorphism
is an isomorphism
\[
\bigwedge^2R^1f_*\OO_Y\xrightarrow{\simeq}R^2f_*\OO_Y.
\]
Every full fibre has structure-sheaf cohomology dimensions
$(h^0,h^1,h^2)=(1,2,1)$, and these two direct images satisfy
base change at every fibre. This proves these two cases of CDP20
Proposition 5.1 by an index and cup-product argument.
\end{proposition}

% Cited input cdp:st:D.1.4; statement or labelled formula only.
\begin{proposition}\hypertarget{tag-01G2}{}\label{cdp:st:D.1.4}


For any $f:Y\to C$ as in \ref{cdp:st:D.1.2} put $E=R^1f_*\OO_Y$.
The coherent Leray sequence gives
\[
0\longrightarrow H^1(C,\OO_C)\xrightarrow{f^*}H^1(Y,\OO_Y)
\xrightarrow{e}H^0(C,E)\longrightarrow0.               \tag{D2}\hypertarget{tag-019U}{}\label{cdp:eq:D2}
\]
At a smooth fibre where cohomological base change holds, the
restriction homomorphism is exactly the composite
\[
H^1(Y,\OO_Y)\xrightarrow{e}H^0(C,E)
\xrightarrow{\operatorname{ev}_c}E\otimes\C(c)
\xrightarrow{\simeq}H^1(F_c,\OO_{F_c}).               \tag{D3}\hypertarget{tag-019V}{}\label{cdp:eq:D3}
\]
By Statement \CompendiumRef{II}{cdp:st:E4}, $E$ is locally free. For $C=\PP^1$, write
$E=\bigoplus_{j=1}^q\OO(a_j)$.
Then the rank of \eqref{cdp:eq:D3} at a general point is exactly
$\#\{j:a_j\geq0\}$. In particular, the nonnegative splitting
degrees claimed by CDP20 Proposition 5.5 are equivalent to its
claimed general-fibre surjectivity when $E$ has rank two and is
locally free.
\end{proposition}

% Cited input cdp:st:D.1.6; statement or labelled formula only.
\begin{proposition}[A second count counterexample]\hypertarget{tag-01G4}{}\label{cdp:st:D.1.6}



\textbf{A second count counterexample under the integral hypotheses.}
Wang's attached construction $X_2(\lambda)$ has integral homology
$H_0=H_6=\Z$, $H_2=H_3=\Z/2\Z$, and the other reduced
homology groups zero. Hence its integral cohomology satisfies
$H^1=H^2=0$. Its fibration has four exceptional values, over which
the reduced fibres are irreducible: the multiplicity-three
bielliptic fibre, two rank-one fillings, and the rank-two filling.
In the rank-one filling over the $I_2$ value, the primitive period
exponent is $(2,1)$; the component count is one, not two.
These are the construction's statements
\texttt{intro:main}, \texttt{fill:rankone},
\texttt{fill:a2}, \texttt{fill:log}, and
\texttt{main:theorem} in the attached primary source
\cite{Wang2026}.

In its cohomology basis the four displayed monodromies are
\[
\begin{aligned}
B_*&=\begin{pmatrix}1&0&-6&-2\\0&-1&-1&1\\0&1&0&-1\\0&0&0&1\end{pmatrix},&\qquad
B_2&=\begin{pmatrix}1&0&6&3\\0&1&2&1\\0&0&1&0\\0&0&0&1\end{pmatrix},\\
B_+&=\begin{pmatrix}1&0&0&0\\0&2&1&0\\0&-1&0&0\\0&0&0&1\end{pmatrix},&\qquad
B_-&=\begin{pmatrix}1&0&0&-1\\0&1&1&0\\0&0&1&0\\0&0&0&1\end{pmatrix}.
\end{aligned}
\]
Their product is one and their common invariant cohomology group
has rank one. On homology the matrices are $B_i^{-t}$; the Smith
diagonal of their difference matrix is $(1,1,1,0)$. Consequently
$t'=1$ and the corrected count is $r=4-1+1=4$, while the printed
count gives seven. The exact matrix calculation is included in the
program. The topological construction and its local fillings are
being used here as cited primary-source results, not inferred from
the matrices alone. Its coherent direct images are derived separately
in Statement \ref{cdp:st:D.1.13} below.
\end{proposition}

% Cited input cdp:st:D.1.9; statement or labelled formula only.
\begin{proposition}\hypertarget{tag-01G7}{}\label{cdp:st:D.1.9}


Let $Y$ be a connected compact complex manifold with
$b_1(Y)=b_2(Y)=0$, and let $f:Y\to\PP^1$ be a proper
holomorphic surjection with connected fibres. There is a class
$\beta\in H^1(Y,\OO_Y)$ such that, for every irrational real number
$r$, the line bundle
\[
L_r=\exp(2\pi i\,r\beta)
\]
has infinite order on every scheme-theoretic fibre of $f$.
Every connected component of $\Pic(Y)$ likewise contains a real
one-parameter family with this property at every irrational real
parameter.
No local-freeness assumption on $R^1f_*\OO_Y$ or constant-dimension
assumption on fibre cohomology is required.
\end{proposition}

% Cited input cdp:st:E1; statement or labelled formula only.
\begin{proposition}[Vertical divisors and the field of meromorphic functions]\hypertarget{tag-01GG}{}\label{cdp:st:E1}


Let $Y$ be a connected compact complex manifold of dimension $n\geq2$
with $H^2(Y,\mathbb R)=0$. Let $g:Y\to C$ be a surjective holomorphic
map to a connected compact Riemann surface, with connected fibres.
Suppose that its smooth fibres over a nonempty open subset of $C$ are
K\"ahler. Every effective divisor of $Y$ is supported on finitely many
fibres, and pullback induces an isomorphism of fields
\[
 g^*: \mathcal M(C)\xrightarrow{\ \sim\ }\mathcal M(Y).
\]
In particular $a(Y)=1$. Reducedness of the other fibres is unnecessary.
\end{proposition}

% Cited input cdp:eq:D3; statement or labelled formula only.
\begin{proposition}\hypertarget{tag-01G2}{}\label{cdp:st:D.1.4}


For any $f:Y\to C$ as in \ref{cdp:st:D.1.2} put $E=R^1f_*\OO_Y$.
The coherent Leray sequence gives
\[
0\longrightarrow H^1(C,\OO_C)\xrightarrow{f^*}H^1(Y,\OO_Y)
\xrightarrow{e}H^0(C,E)\longrightarrow0.               \tag{D2}\hypertarget{tag-019U}{}\label{cdp:eq:D2}
\]
At a smooth fibre where cohomological base change holds, the
restriction homomorphism is exactly the composite
\[
H^1(Y,\OO_Y)\xrightarrow{e}H^0(C,E)
\xrightarrow{\operatorname{ev}_c}E\otimes\C(c)
\xrightarrow{\simeq}H^1(F_c,\OO_{F_c}).               \tag{D3}\hypertarget{tag-019V}{}\label{cdp:eq:D3}
\]
By Statement \CompendiumRef{II}{cdp:st:E4}, $E$ is locally free. For $C=\PP^1$, write
$E=\bigoplus_{j=1}^q\OO(a_j)$.
Then the rank of \eqref{cdp:eq:D3} at a general point is exactly
$\#\{j:a_j\geq0\}$. In particular, the nonnegative splitting
degrees claimed by CDP20 Proposition 5.5 are equivalent to its
claimed general-fibre surjectivity when $E$ has rank two and is
locally free.
\end{proposition}

% Cited input cdp:st:D.1.2; statement or labelled formula only.
\begin{proposition}\hypertarget{tag-01G0}{}\label{cdp:st:D.1.2}


Let $f:Y\to C$ be a proper holomorphic surjection from a connected
compact complex threefold to a compact Riemann surface, with
connected fibres. Let $\Delta$ be a nonempty finite set containing
the critical values, $s=|\Delta|$, $g=g(C)$,
$D=f^{-1}(\Delta)_{\rm red}$, and let $r$ be the number of
irreducible components of $D$. Write $F$ for a smooth fibre and
\[
t'=\dim_{\Q}H_1(F;\Q)_{\pi_1(C\setminus\Delta)}
   =\dim_{\Q}H^1(F;\Q)^{\pi_1(C\setminus\Delta)}.
\]
For $b_2(Y)=0$ the exact component count is
\[
\boxed{\quad r=2g+s-1+t'-b_1(Y).\quad}
\]
Under CDP20's integral assumptions
$H^1(Y;\Z)=H^2(Y;\Z)=0$, the base is $\PP^1$ and
\[
r=s-1+t',\qquad
H_1(F;\Z)_{\pi_1(C\setminus\Delta)}\text{ is free abelian}.
\]
The last assertion concerns the indicated quotient of $H_1(F;\Z)$;
it does not assert that $H_1(F;\Z)$ itself is torsion-free.
\end{proposition}

